\documentclass[aspectratio=169]{beamer}
\usetheme{Madrid}
\usecolortheme{beaver}
\usepackage[utf8]{inputenc}
\usepackage[T1]{fontenc}
\usepackage{tikz}
\usepackage{amssymb}
\usepackage{graphicx}
\setbeamertemplate{navigation symbols}{}
\setbeamercolor{title}{fg=white,bg=cyan!60!black}
\setbeamercolor{frametitle}{fg=white,bg=cyan!50!black}
\setbeamercolor{block title}{bg=cyan!15,fg=black}
\setbeamercolor{block body}{bg=cyan!5,fg=black}

\title{G1 — Sprite Start}
\subtitle{Lesson 1: First Sprite}
\author{ByteForge --- Grade 7 --- Game Maker}
\date{}

\begin{document}

\begin{frame}
    \titlepage
\end{frame}

\begin{frame}{Coordinates — every sprite has an (x, y)}
    \Large
    x (how far right), y (how far down) — y grows DOWN, the opposite of a maths graph. Everything you build this year comes back to these two numbers.

    \begin{center}
    \begin{tikzpicture}[scale=0.9]

        \draw[->, thick] (0,0) -- (5,0) node[right] {$x \to$};
        \draw[->, thick] (0,0) -- (0,-3) node[below] {$y \downarrow$};
        \filldraw[cyan!70!black] (3,-2) circle (2.5pt);
        \node[right] at (3,-2) {(x, y)};
        \node[left] at (-0.1,0.1) {(0,0)};
    \end{tikzpicture}
    \end{center}
\end{frame}

\begin{frame}{Coordinates — every sprite has an (x, y) --- Worked Example}
    \Large
    Worked example: a sprite at x=3, y=2 is 3 steps RIGHT of the left edge and 2 steps DOWN from the top edge — never "2 up", because y only ever counts downward on this canvas.
\end{frame}

\begin{frame}{Your Turn}
    \Large
    Open the student map, find \textbf{G1 — Sprite Start}, and complete
    the \textbf{First Sprite} challenge --- everything you just saw, they now do themselves.
\end{frame}

\end{document}
